\documentclass[10pt,a4paper]{amsart}
\usepackage[margin=19mm]{geometry}
\usepackage{amsmath,amssymb,mathtools}
\usepackage[scr=rsfs,cal=euler]{mathalfa}
\usepackage{xfrac}
\usepackage[unicode,hidelinks,bookmarks=false]{hyperref}
\newcommand{\ie}{i.e.\ }
\newcommand{\cf}{cf.\ }
\newcommand{\resp}{respectively\ }
\newcommand{\Subsection}{\subsection}
\providecommand{\nobreakdash}{\nobreak\hbox{-}}
\setlength{\emergencystretch}{2em}
\multlinegap0pt
\usepackage{bm}
\usepackage{bbm}
\usepackage{dsfont}
\usepackage{xspace}
\usepackage{cancel}
\usepackage{mathtools}
\usepackage{amsthm}
\makeatletter
\@ifundefined{theorem}{\newtheorem{theorem}{Theorem}}{}
\@ifundefined{proposition}{\newtheorem{proposition}[theorem]{Proposition}}{}
\makeatother
\title[Discrete restrictions from Laurent monomial systems for mult]{Discrete restrictions from Laurent monomial systems for multiple Dirichlet series}
\author{Shenghao Hua}
\date{Source: arXiv:2507.23477v3 (CC BY 4.0)}
\begin{document}
\maketitle

\noindent\emph{Source and licence.} Discrete restrictions from Laurent monomial systems for multiple Dirichlet series. Version arXiv:2507.23477v3 (CC BY 4.0), distributed under the Creative Commons Attribution 4.0 International licence, as recorded in the arXiv licence metadata for this exact version and shown on the abstract page https://arxiv.org/abs/2507.23477v3. Full rights notice: licenses/arxiv-2507.23477-CC-BY-4.0.txt.\par\smallskip

\noindent\textit{Editorial scope.} Counted unit: the Proposition whose source label is \texttt{prop:conter} in the article (source0.tex lines 317--319, source label \texttt{prop:conter}), delivered together with the article's own complete original proof of it (source0.tex lines 321--359, one \texttt{proof} environment of 39 lines). No mathematics is written, repaired or re-derived: statement and proof are the article's own text line for line. Required context retained uncounted: none. Every definition, display and label of those lines is retained unchanged. Omitted from this package and not redistributed: 1--316, 320--320, 360--377. Nothing inside a retained excerpt is omitted or altered: every retained body is delivered exactly as the article's lines are written, so no omission receipt of any kind accompanies this package. Citations: the retained excerpts carry no citation command at all, so no bibliography entry of the article is transcribed and no reference is redistributed. Reference adaptation: none was needed and none was made; every reference inside the delivered unit resolves inside it, so no unresolved reference or citation is produced. Printed slips kept literal: no printed word, symbol, hypothesis or number of the counted statement or of its proof was altered. Layout and class: the article's own class is \texttt{amsart} with packages including amsmath, amsthm, amssymb, amsfonts, amscd, mathrsfs, bbm, booktabs, float, appendix, makecell, bbding, hyperref, geometry; the wrapper is the lane's pinned \texttt{amsart} editorial wrapper at 10pt on A4, which restates the theorem-like environments the retained text needs and adds the article's own macro definitions that the retained text needs (none), with extra wrapper packages (bm, bbm, dsfont, xspace, cancel, mathtools). The delivered numbering is the unit's own. Assets: no retained span contains a figure environment, a graphics-inclusion command, a PDF-page inclusion, a table or a TikZ picture; no image file is copied into this package, the package declares no asset dependency and no image file is redistributed with it. Licence: version arXiv:2507.23477v3 of this article is distributed under the Creative Commons Attribution 4.0 International licence, as recorded in the arXiv licence metadata for this exact version and shown on the abstract page https://arxiv.org/abs/2507.23477v3. A full rights notice is delivered at licenses/arxiv-2507.23477-CC-BY-4.0.txt. Characters: every retained excerpt is pure ASCII, so the delivered engine, XeLaTeX with its default Latin Modern text font, reports no missing character.
\begin{proposition}\label{prop:conter}
There exists an irreducible variety \(V\) with infinitely many positive integral points such that \(V\) satisfies Property~(S), but \(V(\mathbb{Z}_{>0})\) is not the positive integral solution set of any Laurent monomial system.
\end{proposition}

\begin{proof}
Let \(V\subset \mathbb{A}^{3}_{\mathbb{Q}}\), with coordinates \((x,y,z)\), be defined by
\begin{equation}\label{eqn:counterexample}
10(x+y-5)^2+(x-y)^2=9,
\end{equation}
where \(z\) is free. After the affine change of variables
\[
u=x+y-5,\qquad v=x-y,
\]
the defining polynomial becomes \(10u^2+v^2-9\). Its homogenization
\[
10u^2+v^2-9w^2
\]
is a nondegenerate ternary quadratic form, so the corresponding plane conic is geometrically irreducible. It follows that \(V\) is irreducible.

If \(x,y\) are positive integers satisfying~\eqref{eqn:counterexample} and \(x+y-5\neq 0\), then the first term on the left-hand side is at least \(10\), which is impossible. Hence \(x+y=5\), and then \((x-y)^2=9\). Therefore
\begin{equation}\label{eqn:counterexample-points}
V(\mathbb{Z}_{>0})
=\{(1,4,n):n\geq 1\}\mathbin{\cup}\{(4,1,n):n\geq 1\}.
\end{equation}

We verify Property~(S). For every prime \(p\neq 2\), the \(p\)-adic valuation vector of the first two coordinates of either family in~\eqref{eqn:counterexample-points} is \((0,0)\). At \(p=2\), it is either \((0,2)\) or \((2,0)\). Thus any prime-by-prime recombination of two points in~\eqref{eqn:counterexample-points} again has first two coordinates equal to either \((1,4)\) or \((4,1)\); its third coordinate is a positive integer. The recombined point therefore still lies in \(V(\mathbb{Z}_{>0})\).

It remains to show that this positive integral solution set cannot arise from a Laurent monomial system. Suppose, to the contrary, that such a system exists. Write one of its equations as
\[
\omega_i x^{a_i}y^{b_i}z^{c_i}=\omega'_i,
\]
where \(a_i,b_i,c_i\in\mathbb{Z}\) and \(\omega_i,\omega'_i\in\mathbb{Z}_{\geq 1}\). Since the equation holds at \((1,4,n)\) for every \(n\geq 1\), comparison of \(n=1\) and \(n=2\) gives \(c_i=0\). Evaluating it at \((1,4,1)\) and \((4,1,1)\) then gives
\[
4^{b_i}=4^{a_i},
\]
and hence \(a_i=b_i\). Consequently, for every \(n\geq 1\),
\[
\omega_i 2^{a_i}2^{b_i}
=\omega_i4^{a_i}
=\omega'_i,
\]
so \((2,2,n)\) satisfies every equation in the system. On the other hand, substituting \((2,2,n)\) into the left-hand side of~\eqref{eqn:counterexample} gives \(10\), not \(9\), a contradiction.
\end{proof}

\end{document}
